\begin{abstract}
The concept of impedance through porous media holds significant importance in acoustical engineering. 
One important problem is optimizing airflow while minimizing noise. 
This paper explores the problem through the study of the acoustic properties of porous rigid media. 
The study examines the acoustic impedance through media with varying pore properties, focusing on the influence of varying pore numbers, percent porosity, and pore length on acoustic impedance. 
By maintaining the material as a constant, the scenario was simplified, enabling the identification of the relationship between pore properties and acoustic impedance. 
Previous studies done on porous media have mostly focused on how the material of the media affects the impedance. 
The results of this study expands the work done by previous studies, offering a new perspective on acoustic impedance. 
The porous media were created using Python to ensure the pores are uniformly spread across the cylinders. 
This experiment uses a traditional impedance tube with a speaker at the top of the tube and two microphones above the porous medium at the bottom. 
We will use frequencies ranging from one hundred fifty Hertz to one thousand Hertz. 
Analysis of reflected waves off the medium will provide the necessary impedance values.  
\end{abstract}
